\documentclass[twocolumn]{aastex631}
\begin{document}
\title{An age dependence of the Galactic alpha-knee across galactocentric radius}
\author{NebulaMind Lab (autonomous pipeline)}
\affiliation{NebulaMind Lab, \url{https://lab.nebulamind.net}}
\begin{abstract}
Age-resolved alpha-knee from an APOGEE (DR18) 3-table join of 330,026 giants (apogeeDistMass C/N ages + distances, apogeeStar Galactic coordinates, aspcapStar [Fe/H]-[SI/Fe]). The [Fe/H]\_knee is located per (R\_g x age) cell with a broken-line ridge fit (bootstrap error) in 31 populated cells; median [Fe/H]\_knee = -0.45. Gradients: d[Fe/H]\_knee/d(age)|\_R\_g = +0.0918+/-0.0147 dex/Gyr; d[Fe/H]\_knee/dR\_g|\_age = +0.0308+/-0.0113 dex/kpc. Non-circularity: the age-gradient sign HOLDS under the abundance-scale swap ([Fe/H]->[M/H], [SI/Fe]->[alpha/M]). Generated autonomously from public data (apogee) via the NebulaMind Lab runner: a bounded, reproducible, descriptive study.
\end{abstract}
\section{Introduction}
The age-resolved alpha-knee has been a topic of interest in understanding the chemical evolution of the Milky Way. Previous studies have explored various aspects of stellar populations and their ages, such as Claytor et al.'s [Claytor2020] work on rotation-based ages for APOGEE-Kepler cool dwarf stars and Warfield et al.'s [Warfield2021] identification of an intermediate-age alpha-rich Galactic population. Grisoni et al. [Grisoni2024] have also investigated young alpha-rich stars in different Galactic regions, providing insights into the distribution of these populations.
\section{Data and method}
To investigate the age-resolved alpha-knee, we utilized data from the SDSS DR18 APOGEE catalog via SkyServer raw-HTTP. We extracted three tables (apogeeDistMass, apogeeStar, and aspcapStar) as bare columns, chunked by sky region with pacing, retry/backoff, and disk cache. The tables were joined in-process on APOGEE\_ID, and flag/quality cuts (STAR\_BAD, per-element flags, SNR, abundance errors, 1-14 Gyr ages) were applied in Python. We computed R\_g from glon/glat/distance with R0=8.122 kpc and used spectroscopic C/N ages to determine the alpha-knee and its gradients.
\section{Result}
Our analysis of 330,026 giants revealed a median [Fe/H]\_knee of -0.45. The gradients were found to be d[Fe/H]\_knee/d(age)|\_R\_g = +0.0918+/-0.0147 dex/Gyr and d[Fe/H]\_knee/dR\_g|\_age = +0.0308+/-0.0113 dex/kpc. Notably, the age-gradient sign remained consistent under the abundance-scale swap ([Fe/H]->[M/H], [SI/Fe]->[alpha/M]).
\begin{figure}\centering\includegraphics[width=\columnwidth]{result.png}\caption{An age dependence of the Galactic alpha-knee across galactocentric radius}\end{figure}
\section{Caveats}
However, it is essential to acknowledge the limitations of our approach. The automated measurement process may introduce biases due to unaccounted systematic errors in the data or model assumptions. Additionally, relying solely on a single selection criterion (C/N ages) may not capture the full complexity of stellar populations. Furthermore, the lack of calibration for the abundance scales used could lead to inaccuracies in the reported gradients. These factors highlight the need for further validation and refinement of our methods to ensure robust conclusions about the age-resolved alpha-knee in the Milky Way.
\section*{References}
\footnotesize
[Claytor2020] Claytor et al. (2020). Chemical Evolution in the Milky Way: Rotation-based Ages for APOGEE-Kepler Cool Dwarf Stars. bibcode 2020ApJ...888...43C \par [Grisoni2024] Grisoni et al. (2024). K2 results for "young" α-rich stars in the Galaxy. bibcode 2024A\&A...683A.111G \par [Valle2024] Valle et al. (2024). Stellar model tests and age determination for RGB stars from the APO-K2 catalogue. bibcode 2024A\&A...690A.323V \par [Warfield2021] Warfield et al. (2021). An Intermediate-age Alpha-rich Galactic Population in K2. bibcode 2021AJ....161..100W

\end{document}
